\section{Evaluation Metrics}\label{sec:metrics}

The s1 paper evaluates test-time scaling with three metrics: control, scaling, and performance \citep{muennighoff2025s1}. The repository mirrors this structure in \texttt{experiments/budget\_forcing/metrics.py}.

\subsection{Control}

Control measures whether the method can reliably reach a target compute budget. In the local implementation, control is computed as the fraction of runs whose observed token count lies within a tolerance band around the target budget. This metric matters because a method that occasionally spends more compute is not sufficient; test-time scaling should be deliberate rather than accidental.

In the current repository state, control is only fully meaningful when target budgets and actual thinking-token counts are both logged in a way that can be matched run by run. For the default \texttt{n\_wait} sweep, the code explicitly notes that control may be unavailable or only partially approximated.

\subsection{Scaling}

Scaling captures whether additional inference-time compute translates into higher accuracy. The local implementation estimates this using the slope of a linear fit over accuracy versus compute level, with accuracy converted to percentage points. Positive scaling indicates that larger compute budgets improve performance; negative scaling suggests that longer reasoning traces are harming the model.

This metric is especially important for a reproduction study because it separates two distinct outcomes. A model may achieve a reasonable peak accuracy while still failing to scale with extra compute, or it may scale positively but only over a narrow range of budgets. In both cases, reporting a single maximum accuracy would obscure the real behavior.

\subsection{Performance}

Performance is defined as the maximum accuracy observed across tested compute levels. It is the simplest metric to interpret and the easiest one to compare against prior work. However, it is also the least informative in isolation. A reproduction that reaches a similar peak performance but only at a much higher compute cost, or with unstable trigger behavior, should not be treated as equivalent to the original claim.

Together, these three metrics provide a balanced view of test-time scaling. Control evaluates budget precision, scaling evaluates the direction and efficiency of extra compute, and performance evaluates the best-achieved outcome.
\label{sec:metrics}

The s1 paper emphasizes that test-time scaling should not be judged by accuracy alone. A useful control method must also manipulate compute in a predictable way. The repository mirrors this framing through three metrics implemented in \texttt{experiments/budget\_forcing/metrics.py}.

\subsection{Control}

Control measures whether the realized compute stays close to the intended budget. In the local implementation, this is defined as the fraction of runs whose actual token count falls within a tolerance band around the target budget. The motivation is straightforward: a method that claims to spend more or less test-time compute must actually do so reliably. If the model cannot follow the intended budget, downstream comparisons become hard to interpret.

\subsection{Scaling}

Scaling measures how accuracy changes as more compute is allocated. The repository uses the slope of an ordinary least squares fit over compute levels and accuracies, scaled into percentage units. Positive scaling indicates that additional compute is associated with better performance; negative scaling suggests wasted or counterproductive computation. This metric is especially relevant for budget forcing because the technique is useful only if longer controlled reasoning often helps instead of merely prolonging incorrect traces.

\subsection{Performance}

Performance is the maximum accuracy achieved across the tested compute levels. This metric captures the best outcome obtained from the current compute-control schedule, regardless of whether all intermediate points are monotonic. It is a practical summary for a coursework reproduction because it answers the simplest applied question: what is the strongest accuracy the current setup reached?

\subsection{Interpretation in This Project}

In the present repository state, control may remain partially unavailable until the run configuration explicitly records target budgets and actual reasoning-token counts for each sample. By contrast, scaling and performance can already be computed from accuracy values across different \texttt{n\_wait} settings. This asymmetry is important and should be reported honestly in the experiments section rather than hidden. A partial metric profile is still informative if the missing piece is clearly explained.